% Môn: Vật lý
% Lớp: 12
% Chương: Chương I. Vật lí nhiệt
% Bài: L12C1 Bài 4. Nhiệt dung riêng · Bài toán công suất, hiệu suất và thời gian đun
% ===== Câu 26 =====
% ID: L12C1B4-88-TL
% Mức: NB
\dangbt{Bài toán công suất, hiệu suất và thời gian đun}
\begin{bt}
% ID: L12C1B4-88-TL
% Mức: NB
Trình bày cách nhận biết và phương pháp giải dạng “Bài toán công suất, hiệu suất và thời gian đun”. Phân tích nhận định sau và sửa lại nếu sai: “Hiệu suất 80% nghĩa là nhiệt có ích lớn hơn năng lượng cung cấp.”
\loigiai{
Trước hết xác định hiện tượng, trạng thái và các đại lượng đã cho. Nêu định nghĩa hoặc hệ thức phù hợp, đổi về đơn vị SI, áp dụng đúng điều kiện và kiểm tra ý nghĩa kết quả. Nhận định cần đối chiếu: Nhiệt có ích bằng H.P.t. Vì vậy nhận định đã cho sai; phát biểu đúng là: Nhiệt có ích bằng H.P.t.
}
\end{bt}

% ===== Câu 30 =====
% ID: L12C1B4-89-TL
% Mức: NB
\begin{bt}
% ID: L12C1B4-89-TL
% Mức: NB
Trình bày cách nhận biết và phương pháp giải dạng “Bài toán công suất, hiệu suất và thời gian đun”. Phân tích nhận định sau và sửa lại nếu sai: “Hiệu suất 80% nghĩa là nhiệt có ích lớn hơn năng lượng cung cấp.”
\loigiai{
Trước hết xác định hiện tượng, trạng thái và các đại lượng đã cho. Nêu định nghĩa hoặc hệ thức phù hợp, đổi về đơn vị SI, áp dụng đúng điều kiện và kiểm tra ý nghĩa kết quả. Nhận định cần đối chiếu: Nhiệt có ích bằng H.P.t. Vì vậy nhận định đã cho sai; phát biểu đúng là: Nhiệt có ích bằng H.P.t.
}
\end{bt}

% ===== Câu 34 =====
% ID: L12C1B4-90-TL
% Mức: TH
\begin{bt}
% ID: L12C1B4-90-TL
% Mức: TH
Trình bày cách nhận biết và phương pháp giải dạng “Bài toán công suất, hiệu suất và thời gian đun”. Phân tích nhận định sau và sửa lại nếu sai: “Nhiệt có ích bằng H.P.t.”
\loigiai{
Trước hết xác định hiện tượng, trạng thái và các đại lượng đã cho. Nêu định nghĩa hoặc hệ thức phù hợp, đổi về đơn vị SI, áp dụng đúng điều kiện và kiểm tra ý nghĩa kết quả. Nhận định cần đối chiếu: Nhiệt có ích bằng H.P.t. Vì vậy nhận định đã cho đúng.
}
\end{bt}

% ===== Câu 38 =====
% ID: L12C1B4-91-TL
% Mức: VD
\begin{bt}
% ID: L12C1B4-91-TL
% Mức: VD
Trình bày cách nhận biết và phương pháp giải dạng “Bài toán công suất, hiệu suất và thời gian đun”. Phân tích nhận định sau và sửa lại nếu sai: “Hiệu suất 80% nghĩa là nhiệt có ích lớn hơn năng lượng cung cấp.”
\loigiai{
Trước hết xác định hiện tượng, trạng thái và các đại lượng đã cho. Nêu định nghĩa hoặc hệ thức phù hợp, đổi về đơn vị SI, áp dụng đúng điều kiện và kiểm tra ý nghĩa kết quả. Nhận định cần đối chiếu: Nhiệt có ích bằng H.P.t. Vì vậy nhận định đã cho sai; phát biểu đúng là: Nhiệt có ích bằng H.P.t.
}
\end{bt}

% ===== Câu 41 =====
% ID: L12C1B4-92-DS
% Mức: NB
\begin{ex}
% ID: L12C1B4-92-DS
% Mức: NB
Xét các phát biểu sau về dạng “Bài toán công suất, hiệu suất và thời gian đun” (bộ số 4).
\choiceTF
{\True Nhiệt có ích bằng H.P.t.}
{Hiệu suất 80% nghĩa là nhiệt có ích lớn hơn năng lượng cung cấp.}
{\True Khi giải dạng “Bài toán công suất, hiệu suất và thời gian đun”, cần xác định đúng trạng thái đầu, trạng thái cuối và quy ước dấu hoặc điều kiện áp dụng.}
{Có thể bỏ qua mọi dữ kiện về khối lượng, nhiệt độ và hiệu suất trong mọi bài toán thuộc dạng này.}
\loigiai{
\begin{itemchoice}
\item (Đúng) Nhiệt có ích bằng H.P.t. (Vì): Nhiệt có ích bằng H.P.t. b) S: Hiệu suất 80% nghĩa là nhiệt có ích lớn hơn năng lượng cung cấp. c) Đ: Khi giải dạng “Bài toán công suất, hiệu suất và thời gian đun”, cần xác định đúng trạng thái đầu, trạng thái cuối và quy ước dấu hoặc điều kiện áp dụng. d) S: Có thể bỏ qua mọi dữ kiện về khối lượng, nhiệt độ và hiệu suất trong mọi bài toán thuộc dạng này. Kết luận dựa trên định nghĩa, điều kiện áp dụng và quy ước trong SGK KNTT
\item (Sai) Hiệu suất 80% nghĩa là nhiệt có ích lớn hơn năng lượng cung cấp.
\item (Đúng) Khi giải dạng “Bài toán công suất, hiệu suất và thời gian đun”, cần xác định đúng trạng thái đầu, trạng thái cuối và quy ước dấu hoặc điều kiện áp dụng.
\item (Sai) Có thể bỏ qua mọi dữ kiện về khối lượng, nhiệt độ và hiệu suất trong mọi bài toán thuộc dạng này.
\end{itemchoice}
}
\end{ex}

% ===== Câu 42 =====
% ID: L12C1B4-93-TL
% Mức: VDC
\begin{bt}
% ID: L12C1B4-93-TL
% Mức: VDC
Trình bày cách nhận biết và phương pháp giải dạng “Bài toán công suất, hiệu suất và thời gian đun”. Phân tích nhận định sau và sửa lại nếu sai: “Nhiệt có ích bằng H.P.t.”
\loigiai{
Trước hết xác định hiện tượng, trạng thái và các đại lượng đã cho. Nêu định nghĩa hoặc hệ thức phù hợp, đổi về đơn vị SI, áp dụng đúng điều kiện và kiểm tra ý nghĩa kết quả. Nhận định cần đối chiếu: Nhiệt có ích bằng H.P.t. Vì vậy nhận định đã cho đúng.
}
\end{bt}

% ===== Câu 45 =====
% ID: L12C1B4-94-DS
% Mức: TH
\begin{ex}
% ID: L12C1B4-94-DS
% Mức: TH
Xét các phát biểu sau về dạng “Bài toán công suất, hiệu suất và thời gian đun” (bộ số 1).
\choiceTF
{Hiệu suất 80% nghĩa là nhiệt có ích lớn hơn năng lượng cung cấp.}
{\True Khi giải dạng “Bài toán công suất, hiệu suất và thời gian đun”, cần xác định đúng trạng thái đầu, trạng thái cuối và quy ước dấu hoặc điều kiện áp dụng.}
{Có thể bỏ qua mọi dữ kiện về khối lượng, nhiệt độ và hiệu suất trong mọi bài toán thuộc dạng này.}
{\True Nhiệt có ích bằng H.P.t.}
\loigiai{
\begin{itemchoice}
\item (Sai) Hiệu suất 80% nghĩa là nhiệt có ích lớn hơn năng lượng cung cấp. (Vì): Hiệu suất 80% nghĩa là nhiệt có ích lớn hơn năng lượng cung cấp. b) Đ: Khi giải dạng “Bài toán công suất, hiệu suất và thời gian đun”, cần xác định đúng trạng thái đầu, trạng thái cuối và quy ước dấu hoặc điều kiện áp dụng. c) S: Có thể bỏ qua mọi dữ kiện về khối lượng, nhiệt độ và hiệu suất trong mọi bài toán thuộc dạng này. d) Đ: Nhiệt có ích bằng H.P.t. Kết luận dựa trên định nghĩa, điều kiện áp dụng và quy ước trong SGK KNTT
\item (Đúng) Khi giải dạng “Bài toán công suất, hiệu suất và thời gian đun”, cần xác định đúng trạng thái đầu, trạng thái cuối và quy ước dấu hoặc điều kiện áp dụng.
\item (Sai) Có thể bỏ qua mọi dữ kiện về khối lượng, nhiệt độ và hiệu suất trong mọi bài toán thuộc dạng này.
\item (Đúng) Nhiệt có ích bằng H.P.t.
\end{itemchoice}
}
\end{ex}

% ===== Câu 46 =====
% ID: L12C1B4-95-TL
% Mức: VDC
\begin{bt}
% ID: L12C1B4-95-TL
% Mức: VDC
Trình bày cách nhận biết và phương pháp giải dạng “Bài toán công suất, hiệu suất và thời gian đun”. Phân tích nhận định sau và sửa lại nếu sai: “Nhiệt có ích bằng H.P.t.”
\loigiai{
Trước hết xác định hiện tượng, trạng thái và các đại lượng đã cho. Nêu định nghĩa hoặc hệ thức phù hợp, đổi về đơn vị SI, áp dụng đúng điều kiện và kiểm tra ý nghĩa kết quả. Nhận định cần đối chiếu: Nhiệt có ích bằng H.P.t. Vì vậy nhận định đã cho đúng.
}
\end{bt}

% ===== Câu 49 =====
% ID: L12C1B4-96-DS
% Mức: TH
\begin{ex}
% ID: L12C1B4-96-DS
% Mức: TH
Xét các phát biểu sau về dạng “Bài toán công suất, hiệu suất và thời gian đun” (bộ số 5).
\choiceTF
{Hiệu suất 80% nghĩa là nhiệt có ích lớn hơn năng lượng cung cấp.}
{\True Khi giải dạng “Bài toán công suất, hiệu suất và thời gian đun”, cần xác định đúng trạng thái đầu, trạng thái cuối và quy ước dấu hoặc điều kiện áp dụng.}
{Có thể bỏ qua mọi dữ kiện về khối lượng, nhiệt độ và hiệu suất trong mọi bài toán thuộc dạng này.}
{\True Nhiệt có ích bằng H.P.t.}
\loigiai{
\begin{itemchoice}
\item (Sai) Hiệu suất 80% nghĩa là nhiệt có ích lớn hơn năng lượng cung cấp. (Vì): Hiệu suất 80% nghĩa là nhiệt có ích lớn hơn năng lượng cung cấp. b) Đ: Khi giải dạng “Bài toán công suất, hiệu suất và thời gian đun”, cần xác định đúng trạng thái đầu, trạng thái cuối và quy ước dấu hoặc điều kiện áp dụng. c) S: Có thể bỏ qua mọi dữ kiện về khối lượng, nhiệt độ và hiệu suất trong mọi bài toán thuộc dạng này. d) Đ: Nhiệt có ích bằng H.P.t. Kết luận dựa trên định nghĩa, điều kiện áp dụng và quy ước trong SGK KNTT
\item (Đúng) Khi giải dạng “Bài toán công suất, hiệu suất và thời gian đun”, cần xác định đúng trạng thái đầu, trạng thái cuối và quy ước dấu hoặc điều kiện áp dụng.
\item (Sai) Có thể bỏ qua mọi dữ kiện về khối lượng, nhiệt độ và hiệu suất trong mọi bài toán thuộc dạng này.
\item (Đúng) Nhiệt có ích bằng H.P.t.
\end{itemchoice}
}
\end{ex}

% ===== Câu 51 =====
% ID: L12C1B4-97-TLN
% Mức: NB
\begin{ex}
% ID: L12C1B4-97-TLN
% Mức: NB
Cho bốn nhận định: (1) Nhiệt có ích bằng H.P.t. (2) Hiệu suất 80% nghĩa là nhiệt có ích lớn hơn năng lượng cung cấp. (3) Cần kiểm tra đơn vị và điều kiện áp dụng khi xử lí dạng “Bài toán công suất, hiệu suất và thời gian đun”. (4) Có thể thay tùy ý nhiệt độ Celsius cho nhiệt độ tuyệt đối trong mọi công thức. Có bao nhiêu nhận định đúng? Trả lời bằng một số nguyên.
\shortans{2}
\loigiai{
Nhận định (1) và (3) đúng; (2) và (4) sai. Vậy có 2 nhận định đúng.
}
\end{ex}

% ===== Câu 53 =====
% ID: L12C1B4-98-DS
% Mức: VD
\begin{ex}
% ID: L12C1B4-98-DS
% Mức: VD
Xét các phát biểu sau về dạng “Bài toán công suất, hiệu suất và thời gian đun” (bộ số 2).
\choiceTF
{\True Khi giải dạng “Bài toán công suất, hiệu suất và thời gian đun”, cần xác định đúng trạng thái đầu, trạng thái cuối và quy ước dấu hoặc điều kiện áp dụng.}
{Có thể bỏ qua mọi dữ kiện về khối lượng, nhiệt độ và hiệu suất trong mọi bài toán thuộc dạng này.}
{\True Nhiệt có ích bằng H.P.t.}
{Hiệu suất 80% nghĩa là nhiệt có ích lớn hơn năng lượng cung cấp.}
\loigiai{
\begin{itemchoice}
\item (Đúng) Khi giải dạng “Bài toán công suất, hiệu suất và thời gian đun”, cần xác định đúng trạng thái đầu, trạng thái cuối và quy ước dấu hoặc điều kiện áp dụng. (Vì): Khi giải dạng “Bài toán công suất, hiệu suất và thời gian đun”, cần xác định đúng trạng thái đầu, trạng thái cuối và quy ước dấu hoặc điều kiện áp dụng. b) S: Có thể bỏ qua mọi dữ kiện về khối lượng, nhiệt độ và hiệu suất trong mọi bài toán thuộc dạng này. c) Đ: Nhiệt có ích bằng H.P.t. d) S: Hiệu suất 80% nghĩa là nhiệt có ích lớn hơn năng lượng cung cấp. Kết luận dựa trên định nghĩa, điều kiện áp dụng và quy ước trong SGK KNTT
\item (Sai) Có thể bỏ qua mọi dữ kiện về khối lượng, nhiệt độ và hiệu suất trong mọi bài toán thuộc dạng này.
\item (Đúng) Nhiệt có ích bằng H.P.t.
\item (Sai) Hiệu suất 80% nghĩa là nhiệt có ích lớn hơn năng lượng cung cấp.
\end{itemchoice}
}
\end{ex}

% ===== Câu 55 =====
% ID: L12C1B4-99-TLN
% Mức: TH
\begin{ex}
% ID: L12C1B4-99-TLN
% Mức: TH
Cho bốn nhận định: (1) Nhiệt có ích bằng H.P.t. (2) Hiệu suất 80% nghĩa là nhiệt có ích lớn hơn năng lượng cung cấp. (3) Cần kiểm tra đơn vị và điều kiện áp dụng khi xử lí dạng “Bài toán công suất, hiệu suất và thời gian đun”. (4) Có thể thay tùy ý nhiệt độ Celsius cho nhiệt độ tuyệt đối trong mọi công thức. Có bao nhiêu nhận định đúng? Trả lời bằng một số nguyên.
\shortans{2}
\loigiai{
Nhận định (1) và (3) đúng; (2) và (4) sai. Vậy có 2 nhận định đúng.
}
\end{ex}

% ===== Câu 56 =====
% ID: L12C1B4-100-TN
% Mức: NB
\begin{ex}
% ID: L12C1B4-100-TN
% Mức: NB
Câu 1 thuộc dạng “Bài toán công suất, hiệu suất và thời gian đun”. Phát biểu nào sau đây đúng?
\choice
{Hiệu suất 80% nghĩa là nhiệt có ích lớn hơn năng lượng cung cấp.}
{Nhiệt lượng và nhiệt độ là hai đại lượng hoàn toàn giống nhau.}
{Mọi quá trình nhiệt học đều không phụ thuộc điều kiện thực hiện.}
{\True Nhiệt có ích bằng H.P.t.}
\loigiai{
Phát biểu đúng là: Nhiệt có ích bằng H.P.t. Các phương án còn lại trái với mô hình, định nghĩa hoặc hệ thức nhiệt học tương ứng.
}
\end{ex}

% ===== Câu 57 =====
% ID: L12C1B4-101-DS
% Mức: VDC
\begin{ex}
% ID: L12C1B4-101-DS
% Mức: VDC
Xét các phát biểu sau về dạng “Bài toán công suất, hiệu suất và thời gian đun” (bộ số 3).
\choiceTF
{Có thể bỏ qua mọi dữ kiện về khối lượng, nhiệt độ và hiệu suất trong mọi bài toán thuộc dạng này.}
{\True Nhiệt có ích bằng H.P.t.}
{Hiệu suất 80% nghĩa là nhiệt có ích lớn hơn năng lượng cung cấp.}
{\True Khi giải dạng “Bài toán công suất, hiệu suất và thời gian đun”, cần xác định đúng trạng thái đầu, trạng thái cuối và quy ước dấu hoặc điều kiện áp dụng.}
\loigiai{
\begin{itemchoice}
\item (Sai) Có thể bỏ qua mọi dữ kiện về khối lượng, nhiệt độ và hiệu suất trong mọi bài toán thuộc dạng này. (Vì): Có thể bỏ qua mọi dữ kiện về khối lượng, nhiệt độ và hiệu suất trong mọi bài toán thuộc dạng này. b) Đ: Nhiệt có ích bằng H.P.t. c) S: Hiệu suất 80% nghĩa là nhiệt có ích lớn hơn năng lượng cung cấp. d) Đ: Khi giải dạng “Bài toán công suất, hiệu suất và thời gian đun”, cần xác định đúng trạng thái đầu, trạng thái cuối và quy ước dấu hoặc điều kiện áp dụng. Kết luận dựa trên định nghĩa, điều kiện áp dụng và quy ước trong SGK KNTT
\item (Đúng) Nhiệt có ích bằng H.P.t.
\item (Sai) Hiệu suất 80% nghĩa là nhiệt có ích lớn hơn năng lượng cung cấp.
\item (Đúng) Khi giải dạng “Bài toán công suất, hiệu suất và thời gian đun”, cần xác định đúng trạng thái đầu, trạng thái cuối và quy ước dấu hoặc điều kiện áp dụng.
\end{itemchoice}
}
\end{ex}

% ===== Câu 59 =====
% ID: L12C1B4-102-TLN
% Mức: VD
\begin{ex}
% ID: L12C1B4-102-TLN
% Mức: VD
Cho bốn nhận định: (1) Nhiệt có ích bằng H.P.t. (2) Hiệu suất 80% nghĩa là nhiệt có ích lớn hơn năng lượng cung cấp. (3) Cần kiểm tra đơn vị và điều kiện áp dụng khi xử lí dạng “Bài toán công suất, hiệu suất và thời gian đun”. (4) Có thể thay tùy ý nhiệt độ Celsius cho nhiệt độ tuyệt đối trong mọi công thức. Có bao nhiêu nhận định đúng? Trả lời bằng một số nguyên.
\shortans{2}
\loigiai{
Nhận định (1) và (3) đúng; (2) và (4) sai. Vậy có 2 nhận định đúng.
}
\end{ex}

% ===== Câu 60 =====
% ID: L12C1B4-103-TN
% Mức: NB
\begin{ex}
% ID: L12C1B4-103-TN
% Mức: NB
Câu 5 thuộc dạng “Bài toán công suất, hiệu suất và thời gian đun”. Phát biểu nào sau đây đúng?
\choice
{Hiệu suất 80% nghĩa là nhiệt có ích lớn hơn năng lượng cung cấp.}
{Nhiệt lượng và nhiệt độ là hai đại lượng hoàn toàn giống nhau.}
{Mọi quá trình nhiệt học đều không phụ thuộc điều kiện thực hiện.}
{\True Nhiệt có ích bằng H.P.t.}
\loigiai{
Phát biểu đúng là: Nhiệt có ích bằng H.P.t. Các phương án còn lại trái với mô hình, định nghĩa hoặc hệ thức nhiệt học tương ứng.
}
\end{ex}

% ===== Câu 63 =====
% ID: L12C1B4-104-TLN
% Mức: VD
\begin{ex}
% ID: L12C1B4-104-TLN
% Mức: VD
Cho bốn nhận định: (1) Nhiệt có ích bằng H.P.t. (2) Hiệu suất 80% nghĩa là nhiệt có ích lớn hơn năng lượng cung cấp. (3) Cần kiểm tra đơn vị và điều kiện áp dụng khi xử lí dạng “Bài toán công suất, hiệu suất và thời gian đun”. (4) Có thể thay tùy ý nhiệt độ Celsius cho nhiệt độ tuyệt đối trong mọi công thức. Có bao nhiêu nhận định đúng? Trả lời bằng một số nguyên.
\shortans{2}
\loigiai{
Nhận định (1) và (3) đúng; (2) và (4) sai. Vậy có 2 nhận định đúng.
}
\end{ex}

% ===== Câu 64 =====
% ID: L12C1B4-105-TN
% Mức: NB
\begin{ex}
% ID: L12C1B4-105-TN
% Mức: NB
Câu 9 thuộc dạng “Bài toán công suất, hiệu suất và thời gian đun”. Phát biểu nào sau đây đúng?
\choice
{Hiệu suất 80% nghĩa là nhiệt có ích lớn hơn năng lượng cung cấp.}
{Nhiệt lượng và nhiệt độ là hai đại lượng hoàn toàn giống nhau.}
{Mọi quá trình nhiệt học đều không phụ thuộc điều kiện thực hiện.}
{\True Nhiệt có ích bằng H.P.t.}
\loigiai{
Phát biểu đúng là: Nhiệt có ích bằng H.P.t. Các phương án còn lại trái với mô hình, định nghĩa hoặc hệ thức nhiệt học tương ứng.
}
\end{ex}

% ===== Câu 67 =====
% ID: L12C1B4-106-TLN
% Mức: VDC
\begin{ex}
% ID: L12C1B4-106-TLN
% Mức: VDC
Cho bốn nhận định: (1) Nhiệt có ích bằng H.P.t. (2) Hiệu suất 80% nghĩa là nhiệt có ích lớn hơn năng lượng cung cấp. (3) Cần kiểm tra đơn vị và điều kiện áp dụng khi xử lí dạng “Bài toán công suất, hiệu suất và thời gian đun”. (4) Có thể thay tùy ý nhiệt độ Celsius cho nhiệt độ tuyệt đối trong mọi công thức. Có bao nhiêu nhận định đúng? Trả lời bằng một số nguyên.
\shortans{2}
\loigiai{
Nhận định (1) và (3) đúng; (2) và (4) sai. Vậy có 2 nhận định đúng.
}
\end{ex}

% ===== Câu 68 =====
% ID: L12C1B4-107-TN
% Mức: TH
\begin{ex}
% ID: L12C1B4-107-TN
% Mức: TH
Câu 2 thuộc dạng “Bài toán công suất, hiệu suất và thời gian đun”. Phát biểu nào sau đây đúng?
\choice
{Nhiệt lượng và nhiệt độ là hai đại lượng hoàn toàn giống nhau.}
{Mọi quá trình nhiệt học đều không phụ thuộc điều kiện thực hiện.}
{\True Nhiệt có ích bằng H.P.t.}
{Hiệu suất 80% nghĩa là nhiệt có ích lớn hơn năng lượng cung cấp.}
\loigiai{
Phát biểu đúng là: Nhiệt có ích bằng H.P.t. Các phương án còn lại trái với mô hình, định nghĩa hoặc hệ thức nhiệt học tương ứng.
}
\end{ex}

% ===== Câu 71 =====
% ID: L12C1B4-108-TLN
% Mức: VDC
\begin{ex}
% ID: L12C1B4-108-TLN
% Mức: VDC
Cho bốn nhận định: (1) Nhiệt có ích bằng H.P.t. (2) Hiệu suất 80% nghĩa là nhiệt có ích lớn hơn năng lượng cung cấp. (3) Cần kiểm tra đơn vị và điều kiện áp dụng khi xử lí dạng “Bài toán công suất, hiệu suất và thời gian đun”. (4) Có thể thay tùy ý nhiệt độ Celsius cho nhiệt độ tuyệt đối trong mọi công thức. Có bao nhiêu nhận định đúng? Trả lời bằng một số nguyên.
\shortans{2}
\loigiai{
Nhận định (1) và (3) đúng; (2) và (4) sai. Vậy có 2 nhận định đúng.
}
\end{ex}

% ===== Câu 72 =====
% ID: L12C1B4-109-TN
% Mức: TH
\begin{ex}
% ID: L12C1B4-109-TN
% Mức: TH
Câu 6 thuộc dạng “Bài toán công suất, hiệu suất và thời gian đun”. Phát biểu nào sau đây đúng?
\choice
{Nhiệt lượng và nhiệt độ là hai đại lượng hoàn toàn giống nhau.}
{Mọi quá trình nhiệt học đều không phụ thuộc điều kiện thực hiện.}
{\True Nhiệt có ích bằng H.P.t.}
{Hiệu suất 80% nghĩa là nhiệt có ích lớn hơn năng lượng cung cấp.}
\loigiai{
Phát biểu đúng là: Nhiệt có ích bằng H.P.t. Các phương án còn lại trái với mô hình, định nghĩa hoặc hệ thức nhiệt học tương ứng.
}
\end{ex}

% ===== Câu 76 =====
% ID: L12C1B4-110-TN
% Mức: TH
\begin{ex}
% ID: L12C1B4-110-TN
% Mức: TH
Câu 10 thuộc dạng “Bài toán công suất, hiệu suất và thời gian đun”. Phát biểu nào sau đây đúng?
\choice
{Nhiệt lượng và nhiệt độ là hai đại lượng hoàn toàn giống nhau.}
{Mọi quá trình nhiệt học đều không phụ thuộc điều kiện thực hiện.}
{\True Nhiệt có ích bằng H.P.t.}
{Hiệu suất 80% nghĩa là nhiệt có ích lớn hơn năng lượng cung cấp.}
\loigiai{
Phát biểu đúng là: Nhiệt có ích bằng H.P.t. Các phương án còn lại trái với mô hình, định nghĩa hoặc hệ thức nhiệt học tương ứng.
}
\end{ex}

% ===== Câu 80 =====
% ID: L12C1B4-111-TN
% Mức: VD
\begin{ex}
% ID: L12C1B4-111-TN
% Mức: VD
Câu 3 thuộc dạng “Bài toán công suất, hiệu suất và thời gian đun”. Phát biểu nào sau đây đúng?
\choice
{Mọi quá trình nhiệt học đều không phụ thuộc điều kiện thực hiện.}
{\True Nhiệt có ích bằng H.P.t.}
{Hiệu suất 80% nghĩa là nhiệt có ích lớn hơn năng lượng cung cấp.}
{Nhiệt lượng và nhiệt độ là hai đại lượng hoàn toàn giống nhau.}
\loigiai{
Phát biểu đúng là: Nhiệt có ích bằng H.P.t. Các phương án còn lại trái với mô hình, định nghĩa hoặc hệ thức nhiệt học tương ứng.
}
\end{ex}

% ===== Câu 84 =====
% ID: L12C1B4-112-TN
% Mức: VD
\begin{ex}
% ID: L12C1B4-112-TN
% Mức: VD
Câu 7 thuộc dạng “Bài toán công suất, hiệu suất và thời gian đun”. Phát biểu nào sau đây đúng?
\choice
{Mọi quá trình nhiệt học đều không phụ thuộc điều kiện thực hiện.}
{\True Nhiệt có ích bằng H.P.t.}
{Hiệu suất 80% nghĩa là nhiệt có ích lớn hơn năng lượng cung cấp.}
{Nhiệt lượng và nhiệt độ là hai đại lượng hoàn toàn giống nhau.}
\loigiai{
Phát biểu đúng là: Nhiệt có ích bằng H.P.t. Các phương án còn lại trái với mô hình, định nghĩa hoặc hệ thức nhiệt học tương ứng.
}
\end{ex}

% ===== Câu 88 =====
% ID: L12C1B4-113-TN
% Mức: VDC
\begin{ex}
% ID: L12C1B4-113-TN
% Mức: VDC
Câu 4 thuộc dạng “Bài toán công suất, hiệu suất và thời gian đun”. Phát biểu nào sau đây đúng?
\choice
{\True Nhiệt có ích bằng H.P.t.}
{Hiệu suất 80% nghĩa là nhiệt có ích lớn hơn năng lượng cung cấp.}
{Nhiệt lượng và nhiệt độ là hai đại lượng hoàn toàn giống nhau.}
{Mọi quá trình nhiệt học đều không phụ thuộc điều kiện thực hiện.}
\loigiai{
Phát biểu đúng là: Nhiệt có ích bằng H.P.t. Các phương án còn lại trái với mô hình, định nghĩa hoặc hệ thức nhiệt học tương ứng.
}
\end{ex}

% ===== Câu 92 =====
% ID: L12C1B4-114-TN
% Mức: VDC
\begin{ex}
% ID: L12C1B4-114-TN
% Mức: VDC
Câu 8 thuộc dạng “Bài toán công suất, hiệu suất và thời gian đun”. Phát biểu nào sau đây đúng?
\choice
{\True Nhiệt có ích bằng H.P.t.}
{Hiệu suất 80% nghĩa là nhiệt có ích lớn hơn năng lượng cung cấp.}
{Nhiệt lượng và nhiệt độ là hai đại lượng hoàn toàn giống nhau.}
{Mọi quá trình nhiệt học đều không phụ thuộc điều kiện thực hiện.}
\loigiai{
Phát biểu đúng là: Nhiệt có ích bằng H.P.t. Các phương án còn lại trái với mô hình, định nghĩa hoặc hệ thức nhiệt học tương ứng.
}
\end{ex}
